% @file paper71_p_vs_np_de.tex
% @brief Paper 71 (DE): Das P-vs-NP-Problem — Komplexitätsklassen, Barrieren und Implikationen
% @author Michael Fuhrmann
% @address Specialist Mathematics Project, Build 147
% @date 2026-03-12

\documentclass[12pt,a4paper]{amsart}

% -----------------------------------------------------------------------
% Pakete
% -----------------------------------------------------------------------
\usepackage{amsmath,amssymb,amsthm}
\usepackage{mathrsfs}
\usepackage{enumitem}
\usepackage{hyperref}
\usepackage{geometry}
\usepackage{microtype}
\usepackage{booktabs}
\usepackage[ngerman]{babel}
\usepackage{tikz}

\geometry{margin=2.5cm}

% -----------------------------------------------------------------------
% Theorem-Environments (Deutsch)
% -----------------------------------------------------------------------
\newtheorem{theorem}{Satz}[section]
\newtheorem{lemma}[theorem]{Lemma}
\newtheorem{corollary}[theorem]{Korollar}
\newtheorem{proposition}[theorem]{Proposition}
\theoremstyle{definition}
\newtheorem{definition}[theorem]{Definition}
\newtheorem{example}[theorem]{Beispiel}
\newtheorem{remark}[theorem]{Bemerkung}
\newtheorem{conjecture}[theorem]{Vermutung}
\newtheorem{openproblem}[theorem]{Offenes Problem}

% -----------------------------------------------------------------------
% Makros
% -----------------------------------------------------------------------
\newcommand{\Pclass}{\mathsf{P}}
\newcommand{\NP}{\mathsf{NP}}
\newcommand{\coNP}{\mathsf{coNP}}
\newcommand{\PSPACE}{\mathsf{PSPACE}}
\newcommand{\EXPTIME}{\mathsf{EXPTIME}}
\newcommand{\EXPSPACE}{\mathsf{EXPSPACE}}
\newcommand{\Lclass}{\mathsf{L}}
\newcommand{\NL}{\mathsf{NL}}
\newcommand{\PH}{\mathsf{PH}}
\newcommand{\BQP}{\mathsf{BQP}}
\newcommand{\NC}{\mathsf{NC}}
\newcommand{\SAT}{\textsc{Sat}}
\newcommand{\TSAT}{\textsc{3-Sat}}
\newcommand{\CLIQUE}{\textsc{Clique}}
\newcommand{\VC}{\textsc{Vertex-Cover}}
\newcommand{\SUBSET}{\textsc{Subset-Sum}}
\newcommand{\HAM}{\textsc{Ham-Cycle}}
\newcommand{\polyredu}{\leq_{\mathrm{p}}}
\newcommand{\N}{\mathbb{N}}
\newcommand{\Z}{\mathbb{Z}}
\newcommand{\F}{\mathbb{F}}
\DeclareMathOperator{\SIZE}{SIZE}
\DeclareMathOperator{\DTIME}{DTIME}
\DeclareMathOperator{\NTIME}{NTIME}
\DeclareMathOperator{\DSPACE}{DSPACE}
\DeclareMathOperator{\NSPACE}{NSPACE}

% -----------------------------------------------------------------------
% Metadaten
% -----------------------------------------------------------------------
\title[Das P-vs-NP-Problem]{%
  Das P-vs-NP-Problem:\\
  Komplexitätsklassen, Barrieren und Implikationen}

\author{Michael Fuhrmann}
\address{Specialist Mathematics Project, Build 147}
\email{specialist-maths@project.local}
\date{12.\ März 2026}

\keywords{P vs.\ NP, NP-Vollständigkeit, Cook-Levin-Satz, Polynomhierarchie,
  Schaltkreiskomplexität, Natural Proofs, Relativierung, Algebrisierung, PCP-Theorem}

\subjclass[2020]{68Q15, 68Q17, 03D15, 68Q25}

% -----------------------------------------------------------------------
% Abstract (VOR \maketitle)
% -----------------------------------------------------------------------
\begin{abstract}
Die Frage, ob $\Pclass = \NP$ gilt, ist das bedeutendste offene Problem der
theoretischen Informatik und eines der sieben Millennium-Probleme des Clay
Mathematics Institute.  Informell fragt sie: Kann jedes Problem, dessen Lösung
\emph{effizient überprüfbar} ist, auch \emph{effizient gelöst} werden?

Dieser Artikel gibt eine in sich geschlossene Darstellung der zentralen
Konzepte.  Wir führen die deterministischen und nicht-deterministischen
Komplexitätsklassen $\Pclass$, $\NP$, $\PSPACE$ und verwandte Klassen ein
und beweisen den \emph{Satz von Cook und Levin} (1971): \SAT{} ist
$\NP$-vollständig.  Wir behandeln Karps 21 $\NP$-vollständige Probleme (1972),
die Polynomhierarchie $\PH$, Schaltkreiskomplexität sowie die drei
grundlegenden \emph{Beweisbarrieren}---Relativierung (Baker--Gill--Solovay
1975, bewiesen), Algebrisierung (Aaronson--Wigderson 2009, bewiesen) und
Natural Proofs (Razborov--Rudich 1994, bewiesen).  Abschließend diskutieren
wir das PCP-Theorem und seine Konsequenzen für Inapproximierbarkeit sowie
die dramatischen Auswirkungen eines Beweises von $\Pclass = \NP$ auf
Kryptographie, Künstliche Intelligenz, Biologie und Mathematik.

Wir unterscheiden stets sorgfältig zwischen \emph{bewiesenen Sätzen} und
\emph{offenen Vermutungen}: Die zentrale Frage $\Pclass \neq \NP$ bleibt eine
\textbf{Vermutung}.
\end{abstract}

\maketitle

\tableofcontents

% -----------------------------------------------------------------------
\section{Einleitung}
% -----------------------------------------------------------------------

Wenige Fragen in der gesamten Mathematik sind zugleich so leicht zu formulieren
und so hartnäckig ungelöst.  Im Jahr 1971 fragte Stephen Cook, ob es einen
effizienten Algorithmus zur Erfüllbarkeitsprüfung Boolescher Formeln gibt.
Präziser fragte er, ob die Klasse $\NP$ der Probleme, deren Lösungen schnell
\emph{überprüft} werden können, mit der Klasse $\Pclass$ der Probleme
übereinstimmt, die schnell \emph{gelöst} werden können.

\begin{openproblem}[$\Pclass$ vs.\ $\NP$, Cook 1971]
Gilt $\Pclass = \NP$?
\end{openproblem}

Das Clay Mathematics Institute nahm dieses Problem im Jahr 2000 in die Liste
der Millennium-Probleme auf und lobte eine Prämie von einer Million US-Dollar
für eine Lösung in beliebiger Richtung aus~\cite{smale1998}.  Trotz jahrzehntelanger
intensiver Arbeit der besten Logiker, Informatiker und Mathematiker der Welt
bleibt die Frage vollständig offen.

Die Bedeutung des Problems geht weit über die theoretische Informatik hinaus.
Ein Beweis von $\Pclass = \NP$ würde implizieren, dass im Wesentlichen jedes
Problem, für das eine Lösung in Polynomialzeit erkannt werden kann, auch in
Polynomialzeit \emph{gefunden} werden kann.  Dies würde die Grundlagen moderner
Public-Key-Kryptographie erschüttern, Proteinfaltung und Wirkstoffdesign
berechenbar machen und die mathematische Beweisfindung automatisieren.
Die überwiegende Mehrheit der Forscher \emph{vermutet} $\Pclass \neq \NP$,
doch kein Beweis existiert.

Das Papier ist wie folgt gegliedert.  Abschnitt~\ref{sec:klassen} führt die
grundlegenden Komplexitätsklassen ein.  Abschnitt~\ref{sec:npv} behandelt
NP-Vollständigkeit und den Satz von Cook--Levin.
Abschnitt~\ref{sec:pvsnp} formuliert die Frage präzise.
Abschnitt~\ref{sec:raum} behandelt Raumkomplexität.
Abschnitt~\ref{sec:ph} stellt die Polynomhierarchie vor.
Abschnitt~\ref{sec:schaltkreis} diskutiert Schaltkreiskomplexität.
Abschnitt~\ref{sec:barrieren} erläutert die drei Hauptbarrieren.
Abschnitt~\ref{sec:implikationen} untersucht Konsequenzen von $\Pclass = \NP$.
Abschnitt~\ref{sec:pcp} präsentiert das PCP-Theorem und Inapproximierbarkeit.

% -----------------------------------------------------------------------
\section{Komplexitätsklassen}
\label{sec:klassen}
% -----------------------------------------------------------------------

\subsection{Turingmaschinen und Zeitkomplexität}

Eine \emph{deterministische Turingmaschine} (DTM) $M$ besteht aus einer
endlichen Zustandsmenge~$Q$, einem Bandalphabet~$\Gamma$, einer
Übergangsfunktion
$\delta\colon Q\times\Gamma\to Q\times\Gamma\times\{L,R,N\}$,
einem Anfangszustand~$q_0$ und einer Menge von Akzeptier- und
Ablehnzuständen.  Bei Eingabe $x\in\{0,1\}^*$ der Länge $n=|x|$ ist die
\emph{Zeitkomplexität} von $M$ auf~$x$ die Anzahl der Schritte bis zum Halt.

\begin{definition}[Zeitkomplexitätsklasse $\DTIME$]
Für eine Funktion $T\colon\N\to\N$:
\[
  \DTIME(T(n)) \;=\;
  \bigl\{\,L\subseteq\{0,1\}^* \;\big|\;
    \text{eine DTM entscheidet }L\text{ in höchstens }c\cdot T(n)\text{ Schritten}\,\bigr\}.
\]
\end{definition}

\begin{definition}[Die Klasse $\Pclass$]
\[
  \Pclass \;=\; \bigcup_{k\geq 1} \DTIME(n^k).
\]
$\Pclass$ ist die Klasse der Entscheidungsprobleme, die von einer DTM in
Polynomialzeit lösbar sind.  Sie ist robust gegenüber vernünftigen Änderungen
des Maschinenmodells (Mehrbandmaschinen, Registermaschinen usw.)\ und wird
weithin als die formale Klasse der \emph{effizienten Berechenbarkeit} angesehen.
\end{definition}

\subsection{Nicht-deterministische Turingmaschinen und $\NP$}

Eine \emph{nicht-deterministische Turingmaschine} (NTM) hat eine
Übergangsrelation, die in jedem Schritt mehrere Möglichkeiten zulässt.
Eine NTM \emph{akzeptiert} $x$, wenn mindestens ein Berechnungspfad in einem
Akzeptierzustand endet.

\begin{definition}[$\NTIME$ und $\NP$]
\[
  \NTIME(T(n)) \;=\;
  \bigl\{\,L \;\big|\;
    \text{eine NTM entscheidet }L\text{ in höchstens }T(n)\text{ Schritten auf allen Pfaden}\,\bigr\}.
\]
\[
  \NP \;=\; \bigcup_{k\geq 1} \NTIME(n^k).
\]
\end{definition}

Eine äquivalente, oft intuitivere Charakterisierung verwendet
\emph{polynomielle Zertifikate}.

\begin{proposition}[Zertifikat-Charakterisierung von $\NP$]
$L\in\NP$ genau dann, wenn es ein Polynom $p$ und einen deterministischen
polynomialzeitigen Verifizierer $V$ gibt mit:
\[
  x\in L \;\iff\; \exists\,c\in\{0,1\}^{p(|x|)}: V(x,c) = 1.
\]
\end{proposition}

\begin{example}
  \textsc{Graph-3-Färbung}: Gegeben ein Graph $G$, existiert eine korrekte
  3-Färbung der Knoten?  Ein Zertifikat ist die Färbung selbst; ihre
  Überprüfung benötigt $O(|E|)$ Zeit.  Also liegt
  \textsc{Graph-3-Färbung}$\in\NP$ vor.
\end{example}

\subsection{Die Inklusion $\Pclass\subseteq\NP$}

\begin{proposition}
$\Pclass\subseteq\NP$.
\end{proposition}
\begin{proof}
Falls $L\in\Pclass$, dient die DTM, die $L$ entscheidet, als Verifizierer,
der sein Zertifikat ignoriert; die Berechnung benötigt höchstens Polynomialzeit.
\end{proof}

Ob die umgekehrte Inklusion gilt---ob $\Pclass = \NP$---ist die zentrale offene
Frage dieses Papiers.

% -----------------------------------------------------------------------
\section{NP-Vollständigkeit}
\label{sec:npv}
% -----------------------------------------------------------------------

\subsection{Polynomialzeit-Reduktionen}

\begin{definition}[Many-one-Reduktion]
Eine Sprache $A$ ist \emph{polynomialzeit-many-one-reduzierbar} auf~$B$,
geschrieben $A\polyredu B$, wenn es eine in Polynomialzeit berechenbare
Funktion $f\colon\{0,1\}^*\to\{0,1\}^*$ gibt mit $x\in A\iff f(x)\in B$.
\end{definition}

\begin{definition}[$\NP$-Schwere und $\NP$-Vollständigkeit]
$B$ ist \emph{$\NP$-schwer}, wenn $A\polyredu B$ für alle $A\in\NP$ gilt.
$B$ ist \emph{$\NP$-vollständig}, wenn $B\in\NP$ und $B$ $\NP$-schwer ist.
\end{definition}

\subsection{Der Satz von Cook und Levin}

Das Boolesche Erfüllbarkeitsproblem \SAT{} fragt: Gegeben eine
Aussagenlogik-Formel $\varphi$ in KNF (konjunktive Normalform) mit Variablen
$x_1,\ldots,x_n$, existiert eine Belegung, die alle Klauseln erfüllt?

\begin{theorem}[Cook--Levin, 1971~\cite{cook1971}, BEWIESEN]
\label{thm:cooklevin}
\SAT{} ist $\NP$-vollständig.
\end{theorem}

\begin{proof}[Beweisskizze]
\textbf{(\SAT{}$\in\NP$).}  Eine erfüllende Belegung dient als polynomielles
Zertifikat; die Verifikation ist linear in der Formelgröße.

\textbf{(\SAT{} ist $\NP$-schwer).}  Sei $L\in\NP$ mit Verifizierer $V$,
der in Zeit $p(n)$ läuft.  Wir kodieren die Berechnung von $V(x,c)$ als
KNF-Formel $\varphi_x$ der Größe $O(p(n)^2)$, deren Variablen das
\emph{Tableau} der Berechnung von~$V$ darstellen: Zustand, Kopfposition und
Bandinhalt zu jedem Zeitschritt.  Die Tableau-Klauseln erzwingen:
\begin{enumerate}[label=(\roman*)]
  \item die Anfangskonfiguration von $V$ auf Eingabe $x$ und Zertifikat~$c$;
  \item die gültigen Einzelschrittübergänge von~$V$;
  \item dass der Endzustand ein Akzeptierzustand ist.
\end{enumerate}
Die Konstruktion ist polynomiell: $|\varphi_x| = O(p(|x|)^2)$ und in
Polynomialzeit berechenbar.  Man zeigt, dass $\varphi_x$ genau dann erfüllbar
ist, wenn $x\in L$.  Da $L\in\NP$ beliebig war, reduziert jede
$\NP$-Sprache auf \SAT{}.
\end{proof}

\subsection{Karps 21 NP-vollständige Probleme}

Im Jahr 1972 zeigte Karp~\cite{karp1972} durch Kettenreduktionen ausgehend
von \SAT{}, dass 21 konkrete kombinatorische und graphentheoretische Probleme
$\NP$-vollständig sind:

\begin{theorem}[Karp 1972~\cite{karp1972}, BEWIESEN]
Die folgenden Probleme sind alle $\NP$-vollständig:
\begin{enumerate}[label=(\arabic*)]
  \item \textsc{3-Sat} ($\varphi$ in 3-KNF)
  \item \CLIQUE{} (enthält $G$ eine $k$-Clique?)
  \item \VC{} (Knotenüberdeckung der Größe $\leq k$?)
  \item \textsc{Unabhängige Menge}
  \item \SUBSET{} (summiert eine Teilmenge ganzer Zahlen zu~$t$?)
  \item \HAM{} (enthält $G$ einen Hamiltonkreis?)
  \item \textsc{Partition}, \textsc{Bin-Packing}, \textsc{Graph-Färbung}, \ldots
\end{enumerate}
\end{theorem}

Diese Ergebnisse etablierten NP-Vollständigkeit als zentrales Konzept der
kombinatorischen Optimierung.

% -----------------------------------------------------------------------
\section{Die P-vs-NP-Frage}
\label{sec:pvsnp}
% -----------------------------------------------------------------------

\subsection{Formale Aussage}

\begin{conjecture}[$\Pclass\neq\NP$ --- die zentrale Vermutung, OFFEN]
\label{verm:pvsnp}
$\Pclass \neq \NP$, d.\,h.\ es existieren Probleme in $\NP$, die von keinem
deterministischen Polynomialzeit-Algorithmus gelöst werden können.
\end{conjecture}

Dies ist eine \emph{Vermutung}, kein Satz.  Kein Beweis---in irgendeiner
Richtung---existiert.  Die überwiegende Mehrheit der Komplexitätstheoretiker
glaubt an $\Pclass\neq\NP$~\cite{gasarch2012}, doch dieser Glaube beruht auf
Intuition, empirischen Belegen und dem Vorhandensein der drei Barrieren aus
Abschnitt~\ref{sec:barrieren}.

\subsection{Warum die meisten Forscher $\Pclass\neq\NP$ vermuten}

\begin{enumerate}[label=(\alph*)]
  \item \textbf{Empirische Evidenz.}  Tausende von NP-vollständigen Problemen
    wurden über 50 Jahre lang untersucht; für keines davon wurde ein
    Polynomialzeit-Algorithmus gefunden.
  \item \textbf{Barrieren.}  Alle bekannten Beweistechniken sind nachweislich
    unzureichend (Abschnitt~\ref{sec:barrieren}).
  \item \textbf{Philosophisches Argument.}  Das \emph{Verifizieren} einer
    Lösung erscheint fundamental leichter als das \emph{Finden} einer Lösung;
    diese Intuition teilen die meisten Experten.
\end{enumerate}

\subsection{Konsequenzen von $\Pclass = \NP$}

Falls Vermutung~\ref{verm:pvsnp} falsch wäre---also $\Pclass = \NP$ gelten
würde---hätte dies revolutionäre Konsequenzen:

\begin{itemize}
  \item \textbf{Kryptographie.}  RSA, Diffie--Hellman, AES-Schlüsselplanung
    und praktisch alle Public-Key-Kryptosysteme stützen sich auf die
    angenommene Schwierigkeit von Problemen in~$\NP$.  Falls $\Pclass=\NP$,
    würden Polynomialzeit-Algorithmen für Ganzzahlfaktorisierung, diskreten
    Logarithmus und Graphisomorphismus folgen.  Die gesamte Infrastruktur
    sicherer Kommunikation, Bankwesen, digitaler Signaturen und
    Kryptowährungen (Bitcoin, Ethereum) wäre gebrochen.
  \item \textbf{Künstliche Intelligenz.}  Maschinelles Lernen, Planung und
    Scheduling enthalten $\NP$-schwere Teilprobleme; ihre polynomiale Lösung
    würde KI revolutionieren.
  \item \textbf{Mathematischer Beweis.}  Das Prüfen, ob eine mathematische
    Aussage einen Beweis der Länge $\leq k$ hat, ist $\NP$-vollständig.
    Falls $\Pclass=\NP$, wäre das \emph{Finden} kurzer Beweise in $\Pclass$:
    Mathematik würde prinzipiell automatisierbar.
  \item \textbf{Biologie.}  Proteinfaltungs-Optimierung, Sequenzalignment
    und phylogenetische Rekonstruktion haben $\NP$-harte Kerne; polynomiale
    Algorithmen könnten die Medizin revolutionieren.
\end{itemize}

% -----------------------------------------------------------------------
\section{Raumkomplexität}
\label{sec:raum}
% -----------------------------------------------------------------------

\subsection{Grundlegende Raumklassen}

\begin{definition}[Raumkomplexitätsklassen]
\begin{align*}
  \DSPACE(S(n)) &= \{L \mid \text{eine DTM entscheidet }L\text{ mit höchstens }S(n)\text{ Bandzellen}\},\\
  \NSPACE(S(n)) &= \{L \mid \text{eine NTM entscheidet }L\text{ mit höchstens }S(n)\text{ Bandzellen}\}.
\end{align*}
Die wichtigsten Raumklassen sind:
\[
  \Lclass = \DSPACE(\log n),\quad
  \NL = \NSPACE(\log n),\quad
  \PSPACE = \bigcup_{k\geq 1}\DSPACE(n^k).
\]
\end{definition}

\subsection{Bekannte Inklusionen}

\begin{theorem}[Raum-Zeit-Hierarchie, BEWIESEN]
Die folgenden Inklusionen gelten, und alle sind bekanntermaßen echt:
\[
  \Lclass \;\subsetneq\; \PSPACE \;\subsetneq\; \EXPTIME \;\subsetneq\; \EXPSPACE.
\]
\end{theorem}

Die echte Inklusion $\Pclass\subsetneq\EXPTIME$ folgt aus dem
\emph{Zeithierarchiesatz}; $\PSPACE\subsetneq\EXPTIME$ folgt aus dem
\emph{Raumhierarchiesatz}.

\begin{theorem}[Raumhierarchiesatz, BEWIESEN]
Für raumkonstruierbare Funktionen $s_1(n) = o(s_2(n))$:
$\DSPACE(s_1(n)) \subsetneq \DSPACE(s_2(n))$.
Insbesondere gilt $\Lclass \subsetneq \PSPACE$.
\end{theorem}

\subsection{Satz von Savitch}

\begin{theorem}[Savitch 1970~\cite{savitch1970}, BEWIESEN]
Für jede raumkonstruierbare Funktion $f(n)\geq\log n$:
\[
  \NSPACE(f(n)) \;\subseteq\; \DSPACE\!\left(f(n)^2\right).
\]
Insbesondere gilt $\NL\subseteq\DSPACE(\log^2 n)$ und
$\PSPACE = \mathsf{NPSPACE}$.
\end{theorem}

\begin{proof}[Beweisidee]
Savitchs Algorithmus löst das Graph-Erreichbarkeitsproblem (das
$\NSPACE(f)$-vollständig ist) deterministisch: Um zu prüfen, ob Konfiguration
$C_1$ die Konfiguration $C_2$ in $\leq 2^k$ Schritten erreichen kann, wird ein
Mittelpunkt $C_m$ geraten und rekursiv $C_1\to C_m$ sowie $C_m\to C_2$ in
$\leq 2^{k-1}$ Schritten verifiziert.  Die Rekursionstiefe beträgt $O(f(n))$,
und der Gesamtraumverbrauch ist $O(f(n)^2)$.
\end{proof}

\subsection{Die vollständige Inklusionskette}

Zusammenfassend gilt:
\[
  \Lclass \;\subseteq\; \NL \;\subseteq\; \Pclass \;\subseteq\; \NP \;\subseteq\;
  \PSPACE \;\subseteq\; \EXPTIME,
\]
wobei die erste und letzte Inklusion echt sind, der Status aller anderen
(einschließlich $\Pclass$ vs.\ $\NP$) jedoch offen bleibt.

% -----------------------------------------------------------------------
\section{Die Polynomhierarchie}
\label{sec:ph}
% -----------------------------------------------------------------------

Die \emph{Polynomhierarchie} (PH) verallgemeinert $\NP$ durch alternierende
Quantoren.

\begin{definition}[Polynomhierarchie~\cite{stockmeyer1976}]
Setze $\Sigma_0^\Pclass = \Pi_0^\Pclass = \Delta_0^\Pclass = \Pclass$ und
definiere rekursiv:
\begin{align*}
  \Sigma_{k+1}^\Pclass &= \NP^{\Sigma_k^\Pclass},\\
  \Pi_{k+1}^\Pclass &= \coNP^{\Sigma_k^\Pclass},\\
  \Delta_{k+1}^\Pclass &= \Pclass^{\Sigma_k^\Pclass}.
\end{align*}
Die Polynomhierarchie ist $\PH = \bigcup_{k\geq 0}\Sigma_k^\Pclass$.
\end{definition}

\begin{theorem}[PH vs.\ PSPACE, BEWIESEN]
$\PH \;\subseteq\; \PSPACE$.
\end{theorem}

\begin{theorem}[Kollaps der Polynomhierarchie, BEWIESEN]
Falls $\Pclass = \NP$, dann gilt $\PH = \Pclass$ (die Hierarchie kollabiert
auf ihre erste Stufe).  Allgemeiner: falls $\Sigma_k^\Pclass = \Pi_k^\Pclass$
für irgendein~$k$, dann gilt $\PH = \Sigma_k^\Pclass$.
\end{theorem}

\begin{proof}
$\Sigma_1^\Pclass = \NP = \Pclass = \Pi_1^\Pclass$ gilt trivialerweise.
Durch Induktion fügt Orakelzugang zu $\Sigma_k^\Pclass = \Pclass$ keine
über $\Pclass$ hinausgehende Stärke hinzu.
\end{proof}

Die meisten Experten vermuten, dass PH nicht kollabiert, was ein weiteres
indirektes Indiz für $\Pclass\neq\NP$ liefert.

% -----------------------------------------------------------------------
\section{Schaltkreiskomplexität}
\label{sec:schaltkreis}
% -----------------------------------------------------------------------

\subsection{Boolesche Schaltkreise}

Ein \emph{Boolescher Schaltkreis} $C_n$ ist ein gerichteter azyklischer Graph
mit Eingabeknoten ($x_1,\ldots,x_n$ sowie Konstanten $0,1$), Gatterknoten
(beschriftet mit UND, ODER, NICHT) und einem einzelnen Ausgabegatter.
Die \emph{Größe} $|C_n|$ ist die Anzahl der Gatter; die \emph{Tiefe}
$d(C_n)$ ist der längste Eingabe-Ausgabe-Pfad.

\begin{definition}[Schaltkreiskomplexitätsklassen]
\begin{align*}
  \Pclass/\mathrm{poly} &= \bigcup_{k}\SIZE(n^k), \\
  \NC^i &= \bigcup_{k}\text{Schaltkreise der Tiefe }O(\log^i n)\text{ und polynomieller Größe},\\
  \NC &= \bigcup_{i\geq 1}\NC^i.
\end{align*}
\end{definition}

Es ist bekannt, dass $\Pclass\subseteq\Pclass/\mathrm{poly}$ und $\NC\subseteq\Pclass$.

\subsection{Satz von Barrington}

\begin{theorem}[Barrington 1989~\cite{barrington1989}, BEWIESEN]
$\NC^1$ stimmt mit der Klasse der Sprachen überein, die von polynomialgroßen
Verzweigungsprogrammen der Breite~5 erkannt werden.  Insbesondere kann jede
$\NC^1$-Sprache durch ein Permutations-Verzweigungsprogramm der Breite~5
berechnet werden.
\end{theorem}

Dieses überraschende Ergebnis zeigt die Robustheit von $\NC^1$ und hat die
Erforschung flacher Schaltkreise nachhaltig beeinflusst.

\subsection{Barriere der Natural Proofs}

\begin{theorem}[Razborov--Rudich 1994~\cite{razborov1994}, BEWIESEN]
Falls eine Einwegfunktion existiert (was weithin angenommen wird), dann kann
kein \emph{natürlicher Beweis} $\Pclass$ von $\NP$ trennen (oder auch nur
super-polynomielle untere Schranken für Boolesche Schaltkreise gegen
$\Pclass/\mathrm{poly}$ beweisen).
\end{theorem}

Ein \emph{natürlicher Beweis} ist ein kombinatorisches Argument, das
(i) \emph{nützlich} (für Funktionen außerhalb einer kleinen Komplexitätsklasse
gültig), (ii) \emph{konstruktiv} (die Eigenschaft ist effizient berechenbar)
und (iii) \emph{groß} (für einen nicht vernachlässigbaren Anteil aller
Funktionen gültig) ist.  Nahezu alle bekannten Techniken für untere Schranken
sind natürlich, was erklärt, warum keine super-polynomiellen unteren Schranken
für allgemeine Schaltkreise bewiesen wurden.

% -----------------------------------------------------------------------
\section{Beweisbarrieren}
\label{sec:barrieren}
% -----------------------------------------------------------------------

Drei fundamentale Meta-Sätze erklären, warum bestehende Techniken die Frage
$\Pclass$ vs.\ $\NP$ nicht auflösen können.  \emph{Alle drei Barrieren sind
bewiesene Sätze; was offen bleibt, ist die Frage $\Pclass$ vs.\ $\NP$ selbst.}

\subsection{Relativierung (Baker--Gill--Solovay 1975)}

\begin{theorem}[Baker--Gill--Solovay 1975~\cite{baker1975}, BEWIESEN]
Es existieren Orakel $A$ und $B$ mit $\Pclass^A = \NP^A$ und
$\Pclass^B \neq \NP^B$.
\end{theorem}

\begin{proof}[Beweisskizze]
Für Orakel $A$: Wähle $A$ als ein $\PSPACE$-vollständiges Problem; dann gilt
$\Pclass^A = \NP^A = \PSPACE$.

Für Orakel $B$: Konstruiere $B$ durch Diagonalisierung so, dass die unäre
Sprache $U_B = \{1^n \mid \exists x\in B,\,|x|=n\}$ in
$\NP^B\setminus\Pclass^B$ liegt.  In Schritt~$n$ wird $B$ so gewählt, dass
die $n$-te $\Pclass^B$-Maschine auf Eingabe $1^n$ gestört wird, während ein
String der Länge~$n$ in $B$ aufgenommen wird.
\end{proof}

\textbf{Konsequenz.}  Jeder Beweis von $\Pclass\neq\NP$ muss
\emph{nicht-relativierend} sein---er darf nicht einheitlich für alle
Orakelerweiterungen gelten.

\subsection{Algebrisierung (Aaronson--Wigderson 2009)}

\begin{theorem}[Aaronson--Wigderson 2009~\cite{aaronson2009}, BEWIESEN]
Die natürliche \emph{algebraische Erweiterung} des Relativierungs-Rahmens,
genannt \emph{Algebrisierung}, reicht ebenfalls nicht aus, um $\Pclass$ von
$\NP$ zu trennen.  Formal existieren algebraische Orakel $\tilde{A}$ mit
$\Pclass^{\tilde{A}} = \NP^{\tilde{A}}$ und algebraische Orakel $\tilde{B}$
mit $\Pclass^{\tilde{B}} \neq \NP^{\tilde{B}}$.
\end{theorem}

Die Algebrisierung erfasst die Stärke der \emph{Arithmetisierungstechnik},
die den Beweisen $\mathsf{IP}=\PSPACE$ und dem PCP-Theorem zugrunde liegt.
Selbst diese gefeierten Beweise können also nicht direkt in einen Beweis
von $\Pclass\neq\NP$ umgewandelt werden.

\subsection{Natural Proofs}

Die Razborov--Rudich-Barriere (Abschnitt~\ref{sec:schaltkreis}) zeigt, dass
\emph{konstruktive kombinatorische} Techniken für untere Schranken
unzureichend sind---vorausgesetzt, Einwegfunktionen existieren.

\subsection{Zusammenfassung der Barrieren}

\begin{remark}
Zusammen zeigen die drei Barrieren:
\begin{itemize}
  \item Jeder Beweis von $\Pclass\neq\NP$ muss nicht-relativierend,
  \item nicht-algebrisierend und
  \item nicht-natürlich sein (es sei denn, Einwegfunktionen existieren nicht).
\end{itemize}
Derzeit ist keine Beweistechnik bekannt, die alle drei Anforderungen erfüllt.
Dies erklärt, warum das Problem seit über 50 Jahren ungelöst ist.
\end{remark}

% -----------------------------------------------------------------------
\section{Implikationen von P=NP}
\label{sec:implikationen}
% -----------------------------------------------------------------------

\subsection{Kryptographie}

Public-Key-Kryptographie beruht auf der angenommenen Schwierigkeit
zahlentheoretischer Probleme:
\begin{itemize}
  \item \textbf{RSA}: Sicherheit setzt die Schwierigkeit der
    Ganzzahlfaktorisierung voraus.  Faktorisierung liegt in $\NP$;
    falls $\Pclass=\NP$, folgt ein Polynomialzeit-Faktorisierungsalgorithmus.
  \item \textbf{Diffie--Hellman / Elliptische-Kurven-Kryptographie}: Basiert
    auf dem Diskreten-Logarithmus-Problem, das ebenfalls in $\NP$ liegt.
  \item \textbf{Hash-Funktionen}: Urbild- und Kollisionssuche liegen in $\NP$;
    Polynomialzeit-Algorithmen würden SHA-256, SHA-3 und BLAKE brechen.
  \item \textbf{Symmetrische Verschlüsselung}: Schlüsselrekonstruktion aus
    Klartext-Geheimtext-Paaren für AES liegt in $\NP$.
\end{itemize}

Die gesamte Infrastruktur sicherer Kommunikation, des Bankwesens, elektronischer
Signaturen und von Kryptowährungen (Bitcoin, Ethereum) wäre kompromittiert.

\subsection{Künstliche Intelligenz und Optimierung}

\begin{itemize}
  \item \textbf{Maschinelles Lernen}: Training tiefer Netzwerke enthält
    $\NP$-harte Teilprobleme; deren Lösbarkeit würde KI drastisch
    beschleunigen.
  \item \textbf{Constraint Satisfaction}: SAT-Lösung liegt der Hardware-
    Verifikation und Planung zugrunde; polynomiales SAT würde den Chipdesign
    revolutionieren.
  \item \textbf{Routing und Scheduling}: Travelling Salesman, Job-Scheduling
    und Fahrzeugrouting sind $\NP$-vollständig; exakte polynomiale Algorithmen
    würden die globale Logistik optimieren.
\end{itemize}

\subsection{Mathematik und Beweisfindung}

Das Problem, ob eine mathematische Aussage einen Beweis der Länge $\leq k$ hat,
ist $\NP$-vollständig (für gegebenes Formalsystem).  Falls $\Pclass=\NP$, wäre
das \emph{Finden} kurzer Beweise in $\Pclass$: Mathematik würde prinzipiell
automatisierbar.  Gödel bemerkte dies bereits in einem Brief an von Neumann
(1956) und erkannte, dass diese Frage im Wesentlichen äquivalent zur
$\Pclass$-vs-$\NP$-Frage ist.

\subsection{Biologie und Wirkstoffdesign}

Proteinfaltungs-Optimierung, Sequenzalignment und phylogenetische Rekonstruktion
haben $\NP$-harte Kerne.  Polynomiale Algorithmen könnten die
Wirkstoffentdeckung und die Reaktion auf Pandemien revolutionieren.

% -----------------------------------------------------------------------
\section{Das PCP-Theorem und Inapproximierbarkeit}
\label{sec:pcp}
% -----------------------------------------------------------------------

\subsection{Probabilistisch Überprüfbare Beweise}

\begin{definition}[PCP-Verifizierer]
Ein \emph{$(\log n, 1)$-beschränkter PCP-Verifizierer} für eine Sprache~$L$
ist ein polynomialzeitiger randomisierter Algorithmus $V$, der bei Eingabe~$x$
der Länge~$n$:
\begin{enumerate}[label=(\roman*)]
  \item $O(\log n)$ zufällige Bits verwendet,
  \item $O(1)$ Bits einer Beweiszeichenkette $\pi$ liest,
  \item mit Wahrscheinlichkeit 1 akzeptiert, falls $x\in L$ (Vollständigkeit),
  \item mit Wahrscheinlichkeit $\geq 1/2$ für alle $\pi$ ablehnt, falls
    $x\notin L$ (Korrektheit).
\end{enumerate}
\end{definition}

\begin{theorem}[PCP-Theorem, Arora--Safra / Arora et al.\ 1998~\cite{arora1998}, BEWIESEN]
\[
  \NP = \mathrm{PCP}[\log n,\, 1].
\]
Jede $\NP$-Sprache besitzt einen probabilistisch überprüfbaren Beweis mit
logarithmischem Zufallsaufwand und konstantem Abfragekomplexität.
\end{theorem}

\begin{proof}[Beweisidee]
Der Beweis verläuft über \emph{PCP-Komposition}: Ein PCP-Verifizierer für
\SAT{} wird konstruiert, indem die erfüllende Belegung als Polynom kleinen
Grades über einem endlichen Körper kodiert wird (Arithmetisierung), und dann
Konsistenz mithilfe des \emph{Sumcheck}-Protokolls und
\emph{Low-Degree-Testings} überprüft wird.  Der Verifizierer verwendet
$O(\log n)$ zufällige Bits zur Auswahl von Körperelementen und liest $O(1)$
Werte.
\end{proof}

\subsection{Inapproximierbarkeitsresultate}

Das PCP-Theorem hat tiefgreifende Konsequenzen für Approximationsalgorithmen:

\begin{theorem}[MAX-SAT-Inapproximierbarkeit, BEWIESEN]
Falls $\Pclass\neq\NP$, kann \textsc{Max-3-Sat} in Polynomialzeit nicht
besser als mit Faktor $\frac{7}{8}+\varepsilon$ für beliebiges $\varepsilon>0$
approximiert werden.  (Eine zufällige Belegung erreicht das Verhältnis $7/8$ im Erwartungswert,
und H{\aa}stads Ergebnis zeigt, dass dies im Wesentlichen optimal ist.)
\end{theorem}

\begin{theorem}[\CLIQUE-Inapproximierbarkeit, BEWIESEN]
Falls $\Pclass\neq\NP$, kann $\CLIQUE$ in Polynomialzeit nicht mit konstantem
Faktor approximiert werden~\cite{hastad1999}.
\end{theorem}

\begin{remark}
Das PCP-Theorem liefert damit einen \emph{bedingten} Beweis, dass
Approximierbarkeit schwierig ist: Es beweist nicht $\Pclass\neq\NP$ (das würde
das offene Problem lösen), zeigt aber, dass falls $\Pclass\neq\NP$ gilt,
viele Probleme nicht einmal approximiert werden können.
\end{remark}

% -----------------------------------------------------------------------
\section*{Schlussfolgerung}
% -----------------------------------------------------------------------

Das P-vs-NP-Problem ist das tiefste und folgenreichste offene Problem der
theoretischen Informatik.  Der Satz von Cook--Levin (bewiesen, 1971) und Karps
Reduktionskatalog (bewiesen, 1972) haben die Landschaft der NP-Vollständigkeit
etabliert.  Die Hierarchiesätze und der Satz von Savitch (bewiesen) gliedern
die Struktur der Komplexitätsklassen.  Das PCP-Theorem (bewiesen) verbindet
Approximierbarkeit mit NP-Schwere.  Doch die zentrale Frage selbst---Vermutung
\ref{verm:pvsnp}---bleibt offen.

Die drei Barrieren---Relativierung, Algebrisierung und Natural Proofs---zeigen,
dass jede Auflösung fundamental neue mathematische Ideen erfordert.  Fortschritt
könnte aus der algebraischen Geometrie (Mulmuley's \emph{Geometric Complexity
Theory}), aus der Beweiskomplexität oder aus unerwarteten Richtungen kommen.

Bis ein Beweis vorliegt, bleibt $\Pclass\neq\NP$ die Arbeitsannahme von
Kryptographen, Algorithmikern und Komplexitätstheoretikern weltweit.

% -----------------------------------------------------------------------
% Literatur
% -----------------------------------------------------------------------
\begin{thebibliography}{99}

\bibitem{cook1971}
S.~A.\ Cook,
\textit{The complexity of theorem-proving procedures},
Proc.\ 3.\ ACM-Symposium on Theory of Computing (STOC), S.~151--158, 1971.

\bibitem{karp1972}
R.~M.\ Karp,
\textit{Reducibility among combinatorial problems},
in R.~E.\ Miller und J.~W.\ Thatcher (Hrsg.), \textit{Complexity of Computer
Computations}, Plenum Press, S.~85--103, 1972.

\bibitem{baker1975}
T.~Baker, J.~Gill und R.~Solovay,
\textit{Relativizations of the P=?NP question},
SIAM Journal on Computing, 4(4):431--442, 1975.

\bibitem{savitch1970}
W.~J.\ Savitch,
\textit{Relationships between nondeterministic and deterministic tape complexities},
Journal of Computer and System Sciences, 4(2):177--192, 1970.

\bibitem{stockmeyer1976}
L.~J.\ Stockmeyer,
\textit{The polynomial-time hierarchy},
Theoretical Computer Science, 3(1):1--22, 1976.

\bibitem{barrington1989}
D.~A.\ M.\ Barrington,
\textit{Bounded-width polynomial-size branching programs recognize exactly those
languages in NC$^1$},
Journal of Computer and System Sciences, 38(1):150--164, 1989.

\bibitem{razborov1994}
A.~A.\ Razborov und S.~Rudich,
\textit{Natural proofs},
Journal of Computer and System Sciences, 55(1):24--35, 1997.
(Konferenzversion: STOC 1994.)

\bibitem{aaronson2009}
S.~Aaronson und A.~Wigderson,
\textit{Algebrization: A new barrier in complexity theory},
ACM Transactions on Computation Theory, 1(1):2:1--2:54, 2009.

\bibitem{arora1998}
S.~Arora, C.~Lund, R.~Motwani, M.~Sudan und M.~Szegedy,
\textit{Proof verification and the hardness of approximation problems},
Journal of the ACM, 45(3):501--555, 1998.

\bibitem{hastad1999}
J.~H{\aa}stad,
\textit{Clique is hard to approximate within $n^{1-\varepsilon}$},
Acta Mathematica, 182(1):105--142, 1999.

\bibitem{arora-barak}
S.~Arora und B.~Barak,
\textit{Computational Complexity: A Modern Approach},
Cambridge University Press, 2009.

\bibitem{smale1998}
S.~Smale,
\textit{Mathematical problems for the next century},
Mathematical Intelligencer, 20(2):7--15, 1998.

\bibitem{gasarch2012}
W.~I.\ Gasarch,
\textit{The second P=?NP poll},
SIGACT News, 43(2):53--77, 2012.

\end{thebibliography}

\end{document}
